\subsection{Budgeted Data Acquisition and Replay}

The preceding corollaries treat unique examples and updates as independently
specified resources. A more operational question is how to divide one budget
between acquiring fresh examples and replaying the examples already available.
Let one unique example cost $c_{\rm d}>0$, let one SGD update cost
$c_{\rm u}>0$, and impose
\begin{equation}
 c_{\rm d}N+c_{\rm u}T\le B.
 \label{eq:budget-constraint}
\end{equation}

\begin{corollary}[Budget allocation]
\label{cor:budget}
For a sufficiently small constant $c_0$, choose
\begin{align}
N_B&=\left\lfloor c_0\min\left\{
 n,m^s,(\eta B/c_{\rm u})^{s/a},B/c_{\rm d}
 \right\}\right\rfloor,\nonumber\\
T_B&=\left\lfloor\frac{B-c_{\rm d}N_B}{c_{\rm u}}\right\rfloor.
\label{eq:budget-schedule}
\end{align}
Then anchor-masked clipped replay satisfies
\begin{align}
\E\R(\widehat w_{T_B})=O\bigl(&m^{-(b-1)}+n^{-(b-1)/s}
 +(\eta B/c_{\rm u})^{-(b-1)/a}\nonumber\\
&+(B/c_{\rm d})^{-(b-1)/s}+e^{-cm}\bigr).
\label{eq:budget-rate}
\end{align}
The unmasked output has the same power laws up to logarithms.
\end{corollary}

The derivation makes the resource tradeoff transparent. The constant $c_0$
ensures $T_B\asymp B/c_{\rm u}$, and the choice of $N_B$ balances the replay
term $N^\alpha/(\eta T)$ against the coverage term $N^{-\beta}$, using
$\alpha+\beta=a/s$. Thus, when $n$ and $m$ do not bind and $a>s$,
\begin{equation}
 N_B\asymp(\eta B/c_{\rm u})^{s/a},\qquad
 K_B=T_B/N_B\asymp B^{1-s/a}.
 \label{eq:budget-interpretation}
\end{equation}
Only a sublinear number of new examples is required as the budget grows; most
additional budget is eventually spent replaying them until optimization meets
coverage. When $a=s$, the two scales are both linear and fresh data must retain
a constant budget share. The schedule is independent of $b$: target smoothness
changes how much error lies beyond a learned rank, whereas $a$ and $s$ determine
where the resource frontiers cross. Equation~\eqref{eq:budget-rate} is a
frontier-matched guarantee for this benchmark, not a universal claim that the
same cost split is minimax for arbitrary representations or noisy tasks.
